\section{Which inputs can explain the difference}\label{sec:inputs}

% Generated by scripts/make_paper_tables.py from the frozen results.
% Do not edit by hand; the next build overwrites it.
\begin{table}[tbp]
\centering
\caption{Candidate explanations for the gap, one input at a time, per \$1,000 of principal}
\label{tab:inputs}
\small
\begin{tabular}{l r r r r r}
\toprule
Note & Gap to close & +0.05 correlation & Flat vol & Timing & Call cost \\
\midrule
N1 & $-$27.92 & 9.14 & 62.49 & $-$4.04 & 27.99 \\
N2 & $-$48.72 & 7.70 & 70.21 & $-$7.94 & 22.34 \\
N3 & $-$20.84 & 5.58 & 44.50 & $-$4.56 & 40.09 \\
N4 & $-$19.16 & 3.55 & 31.90 & $-$3.64 & 36.38 \\
N5 & $-$11.41 & 3.32 & 6.65 & $-$3.76 & $-$51.99 \\
N6 & $-$13.29 & 3.10 & 26.71 & $-$4.55 & 34.23 \\
\bottomrule
\end{tabular}
\notetext{\textit{Notes:} Each column is the change in model value from moving one input while the rest are held at their calibrated values, on a common set of paths. \textit{+0.05 correlation} raises every pairwise correlation, confirming the investor is long correlation on every note. \textit{Flat vol} replaces the fitted surface with a single at-the-money volatility per leg. It raises every value, and by more than the gap on five of the six, so the shape of the surface is the one input measured here whose effect both runs toward the disclosed value and is large enough to span the difference. \textit{Timing} aligns the surface to the pricing date rather than the snapshot date and lowers every value, widening the gap. \textit{Call cost} is what the issuer's option to redeem is worth to the issuer, from the decomposition in Table~\ref{tab:decomposition}. No single column closes the gap in the first; \S\ref{sec:inputs} takes them in turn.}
\end{table}

% Objective: Rule candidate explanations in or out one at a time. Correlation, smile, timing, proxy bias, exercise suboptimality.

A candidate explanation has to clear two bars. It must move the model value
toward the disclosed figure rather than away from it, and it must be capable of
moving it far enough. Direction alone is not enough and magnitude alone is
meaningless if the sign is wrong. Table~\ref{tab:inputs} reports each input's
measured effect on a common set of paths, with one input moved at a time and
the rest held at their calibrated values. Taken in turn, four candidates fail
one bar or the other and two survive both. A third survivor sits outside that
table, because it is a dimension the model used here does not have.

\subsection{Ruled out by direction}

Two candidates move the wrong way.

The funding rate was settled in \S\ref{sec:results} and is the first a reader
reaches for, because the filings point there. Raising the discount spread
lowers the model value and widens the gap; closing it requires a negative
spread on every note. Nothing about a funding assumption the issuer declines to
publish repairs that, because the sign is wrong before the magnitude is
discussed.

Observation timing fails the same way. The surface is observed two business
days after five of the notes priced, and aligning it back to each pricing date
using the move in the matching volatility index lowers every value further, by
\$3.64 to \$4.56 across the headline notes. The adjustment is real and it
belongs in the model, but it enlarges the quantity to be explained rather than
reducing it.

\subsection{Right direction, too small}

Two candidates push the right way and cannot reach.

The proxied Dow leg, discussed in \S\ref{sec:data}, raises the affected note by
\$2.61 with a standard error of \$0.73 when corrected on its own. That is about
five per cent of that note's gap, and it touches one note of the six. It is a
real bias with a measured size, and the size is not the size of the problem.

Extrapolating the surface past the longest listed expiry matters only where the
extrapolation is long. Carrying the last observed variance slope forward
instead of freezing the smile raises the five-year note by \$12.36, which is 44
per cent of its gap and far from negligible. On the three notes under two years
the same change is worth \$0.22 to \$0.77. The term structure assumption is
therefore a genuine source of uncertainty on the long end of the sample and
close to irrelevant on the short end, while the gap is present throughout. It
cannot be the common explanation, though it plausibly contributes to the
largest case.

\subsection{Bounded by the exercise policy}

Suboptimal exercise was quantified in \S\ref{sec:results} and the bound it
gives is partial. On three notes the disclosed value sits between the
no-call and optimally-called valuations, consistent with an issuer recognising
two to twenty-nine per cent of the call's worth. On the two five-year notes it
sits above the no-call valuation entirely at a 69 basis point funding spread,
so no exercise assumption reaches it. An explanation that works for three notes
and is arithmetically unavailable for two others is not the explanation,
although it may well be part of the answer for the three.

\subsection{What survives}

Two of the measured inputs clear both bars, and a third candidate clears them
on published evidence without being measurable here at all. None of the three
is published in the filings.

\paragraph{Correlation.} The investor is long correlation on every note, so
raising it raises the model value. The measured sensitivity is \$3.10 to \$9.14
per 0.05 of uniform pairwise correlation. Closing each gap on correlation alone
would take a uniform increase of 0.15 to 0.27, which is three to five times the
perturbation actually measured, obtained by extrapolating a local sensitivity
well beyond the range over which it was computed and therefore an estimate of
the order rather than a precise requirement. The direction is right and there is
a reason to expect a difference: a desk selling a worst-of is short correlation
and marks it above a realised estimate, which is the approximation
\S\ref{sec:method} flagged as the method's largest. Whether the difference is
of this size is not something six notes can settle. The candidate is live and
unresolved.

\paragraph{The shape of the volatility surface.} Replacing each leg's fitted
surface with a single at-the-money volatility raises all six notes, by \$6.65
to \$70.21. This comparison is quoted over all six rather than over the
headline five, because what it measures is a sensitivity rather than a level,
and the barrier-coverage objection of \S\ref{sec:data} bears on the excluded
note's value, not on how that value responds to flattening its surface. On five
of the six the move is larger than the gap: a fully flat surface overshoots the
disclosed value by \$12.64 to \$33.93, and falls short on only one note, by
\$4.87. Somewhere between the quoted smile and no smile at all
there is a surface that reproduces each disclosed figure exactly.

The claim this supports is narrow and the narrowness matters. No desk prices a
worst-of on a flat volatility surface, and nothing here suggests any issuer
did. What the flat-vol valuation provides is a bound: it shows the distance
between this paper's values and the disclosed ones lies inside the range that
disagreement about the shape of the surface can generate, and comfortably
inside it on five notes. A structure whose principal trigger sits 30 to 40 per
cent below the initial level, with coupon barriers 25 to 35 per cent below it,
is priced off the wing of the smile rather than its level, and
the wing is where quotes thin out and modelling choices diverge most. Two
careful desks working from the same quotes can differ there by more than the
amounts in Table~\ref{tab:results}.

\paragraph{The dynamics of that surface.} A third candidate does not appear in
Table~\ref{tab:inputs}, because the model used here cannot produce it. Local
volatility reproduces today's surface and fixes how it evolves; stochastic and
stochastic-local volatility models add a vol-of-vol dimension that a European
repricing gate cannot see. \citet{Farkas2022} put the point sharply for
autocallable barrier reverse convertibles: the stochastic parameters of a
Heston local-stochastic volatility model control the autocallable price
\emph{while leaving the fit to European options unaffected}, and the richer
dynamics typically produce higher prices than local volatility. Higher is the
direction that would close the gap measured here.

Two qualifications stop that from being an answer. Their products are
single-asset and these are worst-of on three underlyings, so the direction is
suggestive rather than transferable. And the candidates are not independent:
\citet{DeCol2017} study equity worst-of autocallables under multi-dimensional
stochastic local volatility and find the correlations between the stochastic
variances central to pricing them, so correlation and volatility dynamics
interact in precisely this product class rather than contributing separately.
Setting them out as two columns of a table, as this section necessarily does,
is itself a simplification.

The consequence for \S\ref{sec:validation} is worth stating plainly. That
section establishes the engine returns the vanilla prices it was calibrated to.
\citet{Farkas2022} establish that two models can agree on every one of those
prices and still disagree on an autocallable. The validation gate is silent on
this dimension by construction, not by oversight.

\subsection{Taken together}

Of the inputs this study can measure, none closes the gap on every note while
moving in the right direction except the shape of the volatility surface, and
that one does so by bounding the problem rather than by identifying a choice.
Funding and timing are excluded by sign. The proxy and the extrapolation are
too small or too narrow. Exercise policy accounts for part of three notes and
none of two. Correlation remains, sufficient in principle at a shift this
sample cannot adjudicate. Beyond what is measured, the dynamics of the surface
remain open, and the published direction favours them.

That is a weaker conclusion than naming a cause and a stronger one than
reporting an unexplained residual. All three surviving candidates are inputs
the disclosure does not contain, and one of them, correlation, is absent even
from the lists of principal inputs that two of these issuers print
(\S\ref{sec:prior}). A reader holding the pricing supplement cannot
distinguish among them, which is a statement about what the disclosure makes
checkable rather than about what any issuer did.
